\ticket{12}{Слепой поиск, Hill Climbing, Монте-Карло и эволюционные алгоритмы}

\begin{examplan}
    \item Определить слепой поиск и перечислить основные методы.
    \item Описать Hill Climbing и его проблемы.
    \item Объяснить метод Монте-Карло.
    \item Описать эволюционные алгоритмы и их этапы.
\end{examplan}

\subsection*{Короткая версия для письменного ответа}

\textbf{Слепой поиск} не использует эвристику близости к цели и опирается только на описание задачи: начальное состояние, операторы и целевой тест. К нему относят BFS, DFS, поиск с равномерной стоимостью, поиск с ограничением глубины и итеративное углубление.

\subsection*{Методы слепого поиска}

\begin{description}
    \item[BFS] раскрывает состояния по уровням, использует очередь FIFO. Полон при конечном ветвлении и оптимален при одинаковой стоимости шагов, но требует много памяти: $O(b^d)$.
    \item[Uniform-Cost Search] выбирает узел с минимальной стоимостью пути $g(n)$. Полон и оптимален при положительных стоимостях.
    \item[DFS] идет по одной ветви в глубину, использует стек. Требует мало памяти $O(bm)$, но не полон при бесконечных путях и не оптимален.
    \item[Depth-Limited Search] ограничивает глубину DFS числом $L$.
    \item[Iterative Deepening DFS] многократно запускает ограниченный DFS с глубинами $0,1,2,\ldots$; сочетает полноту BFS и малую память DFS.
\end{description}

\subsection*{Hill Climbing}

\textbf{Восхождение к вершине} --- локальный жадный поиск, который переходит от текущего состояния к лучшему соседнему состоянию по целевой функции.

Шаги:
\begin{enumerate}
    \item выбрать начальное состояние;
    \item рассмотреть соседей;
    \item перейти к соседу с лучшей оценкой;
    \item остановиться, если все соседи не лучше текущего состояния.
\end{enumerate}

Проблемы:
\begin{itemize}
    \item локальные оптимумы;
    \item плато, где много состояний имеют одинаковую оценку;
    \item гребни, где нужен непрямой путь улучшения;
    \item зависимость от начальной точки.
\end{itemize}

Варианты: простое восхождение, наискорейшее восхождение, стохастическое восхождение, восхождение со случайными рестартами.

\subsection*{Метод Монте-Карло}

\textbf{Метод Монте-Карло} --- класс алгоритмов, использующих повторную случайную выборку для приближенного решения задачи.

В поиске и оптимизации он обычно работает так:
\begin{enumerate}
    \item случайно генерируются возможные решения;
    \item каждое решение оценивается целевой функцией;
    \item выбирается лучшее или усредняется результат.
\end{enumerate}

Методы Монте-Карло просты, применимы к большим и непрерывным пространствам, но дают вероятностный результат и часто сходятся медленно. В играх важен Monte Carlo Tree Search, где узлы оцениваются через случайные симуляции.

\subsection*{Эволюционные алгоритмы}

\textbf{Эволюционные алгоритмы} --- стохастические методы оптимизации, имитирующие естественный отбор. Они работают с популяцией возможных решений.

Общая схема:
\begin{enumerate}
    \item создать начальную популяцию;
    \item вычислить приспособленность каждого индивида;
    \item выбрать родителей;
    \item применить скрещивание и мутацию;
    \item сформировать новое поколение;
    \item повторять до критерия остановки.
\end{enumerate}

Основные виды: генетические алгоритмы, эволюционные стратегии, эволюционное программирование, генетическое программирование.

Эволюционные алгоритмы полезны, когда функция сложная, негладкая, шумная, задана как черный ящик или пространство решений слишком велико для полного перебора. Их минусы --- высокая вычислительная цена и чувствительность к параметрам: размеру популяции, вероятности мутаций, стратегии отбора.

\begin{important}
    \item Слепой поиск не знает, насколько состояние близко к цели.
    \item BFS полон и оптимален при равных шагах, но дорог по памяти; DFS экономит память, но может зациклиться.
    \item Hill Climbing быстрый и простой, но застревает в локальных оптимумах.
    \item Монте-Карло использует случайные выборки и дает приближенные решения.
    \item Эволюционные алгоритмы ищут решение популяцией через отбор, скрещивание и мутацию.
\end{important}

